\section{Integral Domain}
\begin{definition}[Zero Divisor]
    Let $R$ be a ring. A non-zero element $a \in R$ is called a \textit{zero divisor} if there is a non-zero element $b \in R$ such that either $ab = 0$ or $ba = 0$.
\end{definition}
\begin{eg}
    $\overline{3}$ is a zero divisor in $\Z/6\Z$ since $\overline{3} \cdot \overline{2} = 0$.
\end{eg}
\begin{definition}[Unit]
    Assume $R$ has an identity $1$. An element $u \in R$ is called \textit{a unit in} $R$ if there is some $v \in R$ such that $uv = vu = 1$.
\end{definition}
The set of units in $R$ is denoted by $R^{\times}$.
\begin{eg}
    $\overline{5}$ is a unit in $\Z/6\Z$ since $\overline{5} \cdot \overline{5} = \overline{1}$.
\end{eg}
\begin{definition}[Integral Domain]
    A commutative ring with identity is called an \textit{integral domain} if it has no zero divisor.
\end{definition}
\begin{prop}
    If $R$ is an integral domain, then for all $a,b\in R$,
    \begin{enumerate}
        \item If $ab = 0$ then $a = 0$ or $b = 0$.
        \item If $ab = ac$ then $a = 0$ or $b = c$.
    \end{enumerate}
\end{prop}
\begin{propproof}
    We will prove $2$.
    \begin{align*}
        & ab = ac\\
        &\implies ab - ac = 0\\
        &\implies a(b-c) = 0\\
        &\implies a = 0 \text{ or } b=c.
    \end{align*}
\end{propproof}
\begin{lemma}
    Any finite integral domain is a field.
\end{lemma}
\begin{propproof}
    Let $R$ be a finite integral domain. Hence, $R$ is commutative and has an identity $1$. We just have to prove $R$ is division ring, that, is for every $a \in R$ there is a $b \in R$ such that $ab = 1$.\\
    For $a \neq 0 \in R$ define $\varphi_a:R \rightarrow R$ such that $\varphi_a(x) = ax$.
    Since $a \neq 0$, $\varphi_a$ is injective. Since $R$ is finite, $|\varphi_a(R)| = |R|$. Hence, $\varphi_a$ is also surjective, so, the map is bijective.\\
    Therefore, for all $a \neq 0$, there is $b$, such that $\varphi_a(b) = 1 \implies ab = 1$.\\
    Hence, $R$ is a field. 
\end{propproof}
\newpage
\section{Ring Homomorphisms}
\begin{definition}[Ring Homomorphism]
    A \textit{ring homomorphism} $\varphi: R \rightarrow S$ is a map satisfying the following:
    \begin{enumerate}
        \item $\varphi(a + b)  = \varphi(a)+ \varphi(b)$ for all $a, b \in R$,
        \item $\varphi(a b)  = \varphi(a)\varphi(b)$ for all $a, b \in R$.
    \end{enumerate}
\end{definition}
\begin{definition}[Isomorphism]
    A bijective ring homomorphism is called an \textit{isomorphism}.
\end{definition}
\begin{eg}[Trivial Homomorphism]
    $\varphi: R \rightarrow S$ such that $\varphi(r) = 0$.
\end{eg}
\begin{eg}[Identity Homomorphism]
    $\varphi: R \rightarrow R$ such that $\varphi(r) = r$.
\end{eg}
\begin{eg}
    $\varphi: \Z \rightarrow \Z/n\Z$ such that $\varphi(x) = x \mod n$ is homomorphism.
\end{eg}
\begin{explanation}
    For all $a, b \in \Z$, 
    \begin{align*}
        \varphi(a + b)  &=(a+b) \text{ mod } n\\
                        &= a \text{ mod } n + b \text{ mod } n\\
                        &= \varphi(a) + \varphi(b).
    \end{align*}
    \begin{align*}
        \varphi(a b)  &= ab \text{ mod } n\\
                      &= (a \text{ mod } n) (b \text{ mod } n)\\
                      &= \varphi(a) \cdot \varphi(b).
    \end{align*}
\end{explanation}
\begin{definition}[Polynomial Ring]
    % Let $R$ be a ring. Define $R[x]$ to be the set of all polynomials in $x$ with coefficients in $R$. In symbols,
    % \[R[x] = \{a_nx^n + a_{n-1}x^{n-1} + ... + _1x + a_0 : a_n, a_{n-1}, ..., a_1, a_0 \in R\}.\]
    % $R[x]$ is a ring called \textit{polynomial ring}.
    Let $R$ be a ring, $x$ be an indeterminate and define $R[x]$ be the set of all formal sums
    \begin{align*}
        R[x] := a_nx^n + a_{n-1}x^{n-1} + \cdots + a_1x + a_0 = \sum a_ix^i.
    \end{align*}
    Given two polynomials $f = \sum a_ix^i$ and $g = \sum b_ix^i$ in $R[x]$, the sum of $f$ and $g$ is defined as
        \[f + g = \sum (a_i + b_i)x^i,\]
    (where we have implicitly assumed that $m \leq n$ and we set $b_i = 0$ for $i > m$).
    And the product of $f$ and $g$ as
    \[fg = \sum_{i}^{}(\sum_{j}^{}a_jb_{i-j})x^i.\]
    With this rule of addition and multiplication, $R[x]$ becomes a ring, called \textit{polynomial ring}. 
\end{definition}
In the definition above, each $a_i \in R$ is called the coefficients of the polynomial; $a_i$ is the coefficient of $x^i$ and the zero element given as the polynomial with zero coefficients.
\begin{eg}
    $\varphi: \Q[x] \rightarrow \Q$ such that $\varphi(P(x)) = P(0)$ is a ring homomorphism.
\end{eg}
\begin{explanation}
    \begin{align*}
        &\varphi[P_1(x) + P_2(x)]\\
        &= P_1(0) + P_2(0)\\
        &= \varphi[P_1(x)] + \varphi[P_2(x)].
    \end{align*}
     \begin{align*}
        &\varphi[P_1 P_2]\\
        &= P_1P_2(0)\\
        &= P_1(0) P_2(0)\\
        &= \varphi[P_1(x)]  \varphi[P_2(x)].
    \end{align*}
\end{explanation}
\section{Kernel and Image of Ring Homomorphism}
\begin{definition}[Kernel of Ring]
    For rings, the \textit{kernel} is the set of all elements in $R$ that get mapped to $0 \in S$. In symbols,
    \[\ker \varphi = \{r \in R: \varphi(r) = 0\}.\]
\end{definition}
\begin{prop}
    \label{def:kernelisideal}
    Let $R$ and $S$ be rings and $\varphi: R \rightarrow S$ a ring homomorphism. Then,
    \begin{enumerate}
        \item The image of $\varphi$ is a subring of $S$.
        \item ker $\varphi$ is a subring of $R$. Furthermore, ker $\varphi$ is closed under multiplication by every element of $R$.
    \end{enumerate}
\end{prop}
\begin{propproof}
    \begin{enumerate}
        \item Let $a, b \in \im \varphi$, then there is $r_1, r_2 \in R$ such that $a = \varphi(r_1)$ and $b = \varphi(r_2)$. Since $\varphi$ is a homomorphism,
        \[a - b = \varphi(r_1) - \varphi(r_2) = \varphi(r_1 - r_2) \in S.\]
        Hence, im $\varphi$ is closed under subtraction.\\
        Again since $\varphi$ is a homomorphism,
        \[ab = \varphi(r_1)  \varphi(r_2) = \varphi(r_1  r_2) \in S.\]
        Therefore, im $\varphi$ is closed under multiplication.\\
        Hence, im $\varphi$ is a subring of $S$.
        \item Let $a, b \in \ker \varphi$, then $\varphi(a) = 0$ and $\varphi(b) = 0$. Since $\varphi$ is a homomorphism,
        \[\varphi(a - b) = \varphi(a) - \varphi(b) = 0 - 0 = 0.\]
        Hence, $(a-b) \in \ker \varphi$, so ker $\varphi$ is closed under subtraction.\\
        Again since $\varphi$ is homomorphism,
        \[\varphi(ab) = \varphi(a)\varphi(b) = 0 \cdot 0 = 0.\]
        Therefore, $ab \in \ker \varphi$, so ker $\varphi$ is closed under multiplication.\\
        Hence, ker $\varphi$ is a subring of $R$.\\
        Now, for the last part, let \( a \in \ker(\varphi) \) and \( r \in R \). Then \( \varphi(a) = 0 \). Since \( \varphi \) is a ring homomorphism, we have:
        \[
        \varphi(ra) = \varphi(r) \cdot \varphi(a) = \varphi(r) \cdot 0= 0,
        \]
        and
        \[
        \varphi(ar) = \varphi(a) \cdot \varphi(r) = 0 \cdot \varphi(r) = 0.
        \]
        Thus, \( ra, ar \in \ker(\varphi) \). Hence, \( \ker(\varphi) \) is closed under multiplication by elements of \( R \).
    \end{enumerate}
\end{propproof}